\chapter{Marco Teórico}
\pagestyle{fancy}

\section{Arquitecturas de Software}

\subsection{Arquitecturas Monolíticas}
El desarrollo de software tradicional frecuentemente se orienta hacia la construcción de aplicaciones monolíticas. En este tipo de arquitectura, todos los componentes del sistema (interfaz de usuario, lógica de negocio y acceso a datos) se empaquetan y despliegan como una única unidad lógica y ejecutable \cite{richardson2018microservices}. Si bien este enfoque simplifica el desarrollo inicial, las pruebas y el despliegue a pequeña escala, tiende a presentar problemas significativos de escalabilidad y mantenimiento a medida que la base de código crece. Un fallo en un componente puede derivar en la caída de toda la aplicación, y actualizar una funcionalidad requiere el redespliegue del sistema entero.

\subsection{Arquitecturas Basadas en Microservicios}
Para mitigar las deficiencias de los monolitos, la industria evolucionó hacia la arquitectura basada en microservicios. Según Fowler y Lewis (2014) \cite{fowler2014microservices}, los microservicios son un enfoque de desarrollo en el que una aplicación única se construye como una suite de pequeños servicios, donde cada uno ejecuta su propio proceso y se comunica mediante mecanismos ligeros (usualmente APIs HTTP/REST). Estos servicios están construidos alrededor de capacidades de negocio, son desplegables de manera independiente y pueden ser escritos en diferentes lenguajes de programación o utilizar diferentes tecnologías de almacenamiento de datos.

\subsection{Comparación entre Arquitecturas Monolíticas y Microservicios}
La principal diferencia radica en el acoplamiento y la unidad de despliegue. Las arquitecturas monolíticas ofrecen simplicidad operativa en etapas tempranas pero sufren de alto acoplamiento. Por el contrario, los microservicios brindan alta cohesión, escalabilidad independiente por servicio, y resiliencia ante fallos aislados. No obstante, introducen una gran complejidad operativa inherente a los sistemas distribuidos, tales como la latencia de red, la tolerancia a fallos, el monitoreo distribuido y la orquestación de despliegues \cite{richardson2018microservices}.

\section{Computación Autonómica}

\subsection{Concepto de Computación Autonómica}
Inspirada en el funcionamiento del sistema nervioso autónomo del cuerpo humano, la Computación Autonómica fue propuesta por IBM en 2001 para lidiar con la creciente e inmanejable complejidad de la gestión de sistemas informáticos \cite{kephart2003vision}. El concepto define sistemas informáticos capaces de autogestionarse en base a políticas de alto nivel establecidas por administradores humanos, ocultando la complejidad subyacente de las operaciones técnicas.

\subsection{Propiedades de los Sistemas Autonómicos}
Un sistema autonómico se rige típicamente por cuatro propiedades fundamentales (conocidas como las propiedades "Self-CHOP") \cite{mendonca2012architecture}:
\begin{itemize}
    \item \textbf{Auto-configuración (Self-configuration):} Capacidad de instalar y adaptar componentes de software automáticamente sin intervención externa.
    \item \textbf{Auto-recuperación (Self-healing):} Habilidad para detectar, diagnosticar y reparar problemas operativos, interrupciones y fallos en el software o hardware.
    \item \textbf{Auto-optimización (Self-optimization):} Monitorización proactiva y ajuste continuo de los recursos para garantizar un rendimiento óptimo.
    \item \textbf{Auto-protección (Self-protection):} Capacidad de anticipar, detectar e identificar ataques maliciosos o fallos en cascada, tomando acciones para salvaguardar el sistema.
\end{itemize}

\section{Modelo MAPE-K}

\subsection{Componentes del Modelo MAPE-K}
El modelo MAPE-K (Monitor, Analyze, Plan, Execute, Knowledge) es el marco arquitectónico de referencia introducido por IBM en 2003 para construir sistemas autonómicos \cite{kephart2003vision}. Divide el proceso de auto-gestión en cuatro fases apoyadas sobre una base de conocimiento compartida:
\begin{itemize}
    \item \textbf{Monitor (Monitoreo):} Recopila información de topología, métricas y logs de los recursos gestionados.
    \item \textbf{Analyze (Análisis):} Procesa la información para entender el estado actual del sistema y determinar si se ha producido una anomalía.
    \item \textbf{Plan (Planificación):} Estructura las acciones necesarias para alcanzar los objetivos o solucionar la falla detectada.
    \item \textbf{Execute (Ejecución):} Aplica el plan a través de actuadores o comandos sobre el sistema objetivo.
    \item \textbf{Knowledge (Conocimiento):} Base de datos o repositorio histórico que nutre y es nutrido por las otras cuatro fases \cite{ibm2006architectural}.
\end{itemize}

\subsection{MAPE-K en Sistemas Distribuidos modernos (AWARE)}
La transición hacia microservicios introdujo nuevos niveles de complejidad, haciendo que la implementación estricta de un único ciclo MAPE-K resultara un cuello de botella. Nuevos enfoques, como el marco AWARE (AIOps-based Workflow for Autonomic Resiliency and Execution), adaptan este modelo integrando capacidades de Inteligencia Artificial directamente en las fases de Análisis y Planificación. Esto permite que el sistema pase de decisiones basadas en reglas estáticas a razonamientos basados en modelos de lenguaje (LLMs) u otros clasificadores de aprendizaje automático \cite{dang2019aiops}.

\section{Sistemas Auto-adaptativos}

\subsection{Definición y Características}
Un sistema auto-adaptativo tiene la capacidad de modificar su comportamiento y estructura en tiempo de ejecución en respuesta a cambios en su entorno u operaciones. Weyns (2019) establece que un sistema verdaderamente auto-adaptativo requiere retroalimentación constante y mecanismos de toma de decisiones no deterministas que reduzcan la carga cognitiva de los operadores humanos \cite{weyns2019software}.

\subsection{Relación con la Computación Autonómica}
Mientras que la computación autonómica abarca la visión completa de sistemas auto-manejables (auto-reparación, auto-optimización, auto-protección y auto-configuración), la auto-adaptación es el mecanismo de ingeniería de software práctico para implementarlo. Herramientas contemporáneas de aseguramiento de calidad (como análisis de cobertura de \textit{JaCoCo} o pruebas end-to-end con \textit{Cypress}) juegan un rol clave, ya que la auto-adaptación requiere que el código alterado mantenga su validez estructural.

\newpage

\section{Observabilidad en Sistemas Distribuidos}

\subsection{Monitoreo}
En arquitecturas monolíticas, el monitoreo a menudo se reducía a verificar si el servidor estaba encendido (uptime) y si el uso de CPU/RAM era normal. Sridharan (2018) \cite{sridharan2018distributed} distingue entre \textit{monitoreo} (el acto de observar el estado general y predecible de un sistema) y la \textit{observabilidad} (la cualidad de un sistema que permite inferir estados internos basándose únicamente en el conocimiento de sus salidas externas). En microservicios, la observabilidad es obligatoria debido a que las fallas rara vez son predecibles y provienen de la interacción compleja entre múltiples componentes pequeños.

\subsection{Logs y Métricas}
Los pilares de la observabilidad se fundamentan en la agregación de telemetría:
\begin{itemize}
    \item \textbf{Métricas:} Valores numéricos agregados a lo largo del tiempo, útiles para disparar alertas cuando se sobrepasan umbrales (ej. latencia p95 o el uso de CPU) \cite{sridharan2018distributed}. Herramientas como \textit{Prometheus} son el estándar en sistemas modernos.
    \item \textbf{Logs (Registros):} Representan eventos discretos y detallados producidos por las aplicaciones. Al centralizarse en sistemas como \textit{Loki}, los logs proporcionan el contexto exacto ("el qué y el porqué") cuando una métrica señala que algo anda mal.
\end{itemize}

\subsection{Trazabilidad (Tracing)}
La trazabilidad distribuida permite el seguimiento del ciclo de vida de una solicitud HTTP única a través de todos los microservicios involucrados en satisfacerla. Aporta un identificador único (Trace ID) que reconstruye la cadena temporal de llamadas de red y tiempos de espera, vital para debugear latencia entre microservicios \cite{sridharan2018distributed}.

\section{Mantenimiento de Software}

\subsection{Tipos de Mantenimiento}
El estándar internacional ISO/IEC/IEEE 14764 \cite{iso14764} establece que el mantenimiento del software ocurre una vez que el sistema ha sido puesto en producción y lo categoriza en cuatro vertientes:
\begin{itemize}
    \item \textbf{Correctivo:} Modificaciones reactivas para corregir fallos o bugs descubiertos.
    \item \textbf{Adaptativo:} Modificaciones para asegurar que el software siga siendo funcional en un entorno tecnológico cambiante (ej. migraciones de base de datos o de S.O.).
    \item \textbf{Perfectivo:} Mejoras en el rendimiento o implementación de nuevas funcionalidades operativas (sin cambiar la lógica de negocio subyacente).
    \item \textbf{Preventivo:} Refactorización o mitigación técnica para prevenir fallos latentes o corregir la deuda técnica antes de que cause incidencias.
\end{itemize}

\subsection{Retos en Sistemas Distribuidos}
Realizar tareas de mantenimiento en arquitecturas basadas en microservicios introduce retos significativos. La falla no suele concentrarse en una sola línea de código, sino que puede deberse a interrupciones de red, fallos en contenedores terceros o latencias asíncronas \cite{richardson2018microservices}. Identificar la raíz del problema (Root Cause Analysis - RCA) consume una porción abrumadora del tiempo del desarrollador, lo cual incrementa el MTTR (Mean Time to Recovery).

\section{Inteligencia Artificial aplicada al Mantenimiento de Software}

\subsection{Detección de Anomalías}
El uso de algoritmos de Machine Learning e IA permite establecer líneas base (baselines) del comportamiento "normal" de un sistema (tasa de errores, latencia habitual, picos estacionales de tráfico). Cuando los patrones reales se desvían de estas líneas base predictivas, la IA puede alertar al equipo antes de que ocurra una interrupción total del servicio \cite{notaro2021aiops}.

\subsection{Apoyo a la Toma de Decisiones}
El advenimiento de los Modelos de Lenguaje de Gran Escala (LLMs) ha revolucionado la gestión de operaciones de TI. Al abstraer métricas matemáticas, archivos de log crudos y topología de servicios en representaciones de texto, los LLMs son capaces de "entender" y sintetizar errores complejos en diagnósticos legibles por humanos. Actúan como asistentes (Copilots) para los ingenieros DevOps, sugiriendo fragmentos de código o configuraciones óptimas para recuperar o escalar la infraestructura.

\section{AIOps y Automatización de Operaciones}

\subsection{Concepto de AIOps}
El término AIOps (Artificial Intelligence for IT Operations) fue acuñado por Gartner en 2016 \cite{gartner2017aiops} para describir la aplicación de Big Data, Machine Learning y otras tecnologías avanzadas a los flujos de operaciones de TI. AIOps mejora el monitoreo tradicional al correlacionar eventos ruidosos, encontrar la causa raíz (RCA) con mayor velocidad y proporcionar remediación predictiva.

\subsection{Aplicaciones en Sistemas Modernos}
En ecosistemas orquestados por herramientas como Docker o Kubernetes, AIOps puede predecir cuándo un contenedor fallará analizando un aumento anormal (drift) de la memoria o latencia. La integración de AIOps en modelos autonómicos (como MAPE-K) representa el estado del arte en la ingeniería de resiliencia (Resilience Engineering) \cite{notaro2021aiops}.

\section{Integración de IA en Arquitecturas de Microservicios}

\subsection{Desafíos de Integración}
La principal barrera para incorporar Inteligencia Artificial en operaciones es la fragmentación de los datos. En un entorno de microservicios, los registros están distribuidos en múltiples nodos. Para que un LLM o un modelo estadístico funcione, se requiere la implementación de patrones de agregación mediante herramientas como Loki y Prometheus, garantizando un punto de entrada centralizado de contexto \cite{newman2015building}.

\subsection{Oportunidades y Beneficios (Mantenimiento Proactivo)}
Al abstraer las métricas como contexto ("prompts") para modelos generativos, los administradores obtienen diagnósticos en lenguaje natural. Esto democratiza la administración de servidores complejos y acelera drásticamente los tiempos de recuperación (MTTR). Además, permite establecer un canal de auto-mantenimiento donde el sistema puede inferir su propio nivel de riesgo, recomendando a los equipos humanos refactorizar secciones de código o ajustar límites de contenedores antes de que se convierta en una caída masiva.